\chapter{Resultados e Discussão}
\label{cap:resultados}

<Apresente as formas de onda e as medidas extraídas do \texttt{.raw} e do
\texttt{.log}. Cada plot gerado com \texttt{./claw-spice raw plot ... --png}
entra como figura a partir de \texttt{imagens/generated/}.>

\begin{figure}[htb]
  \centering
  % Substitua pelo PNG do plot em imagens/generated/
  %\includegraphics[width=\textwidth]{<plot>.png}
  \fbox{\parbox[c][4cm][c]{\textwidth}{\centering <forma de onda:
  imagens/generated/<plot>.png>}}
  \caption{<Descrição da forma de onda, ex.: resposta ao degrau de
  $V(\mathrm{out})$>.}
  \label{fig:resposta}
\end{figure}

\section{Medidas}
\label{sec:medidas}

<Transcreva os resultados \texttt{.meas} do \texttt{.log}
(\texttt{./claw-spice log summary}) e compare com os valores teóricos.>

\begin{table}[htb]
  \centering
  \caption{Medidas da simulação e valores teóricos.}
  \label{tab:medidas}
  \begin{tabular}{lSSS}
    \toprule
    Grandeza & {Simulado} & {Teórico} & {Erro (\%)} \\
    \midrule
    <constante de tempo $\tau$> & 1.002 & 1.000 & 0.2 \\
    \bottomrule
  \end{tabular}
\end{table}

\section{Discussão}
\label{sec:discussao}

<Discuta as diferenças entre simulado e teórico, avisos relevantes do
\texttt{.log}, limitações do modelo e fontes de erro.>
